\chapter{Poisson processes, intensity, and compensators}
\label{ch:poisson}

\begin{motivation}
Everything we build rests on one fact: a point process whose ``rate forecast'' is
fixed in advance --- not reacting to history --- is a Poisson process, and Poisson
counts are independent and easy. This chapter makes ``rate forecast'' precise
(the \emph{compensator}) and proves the fact (Watanabe). Keep this chapter close;
we invoke it constantly.
\end{motivation}

\section{Counting processes and conditional intensity}

A simple point process produces times $t_1<t_2<\cdots$; its counting process is
$N(t)=\#\{i:t_i\le t\}$, and its internal history is
$\Hist_t=\sigma\{N(s):s\le t\}$. The \emph{conditional intensity} is the
instantaneous event rate given the past,
\begin{equation}
\lambda(t)=\lim_{h\downarrow0}\tfrac1h\,\E\!\big[N(t{+}h)-N(t)\mid\Hist_{t^-}\big],
\end{equation}
so heuristically $\E[\dd N(t)\mid\Hist_{t^-}]=\lambda(t)\dd t$.

\section{The compensator}

By the Doob--Meyer decomposition there is a unique predictable increasing process
$\Lambda$ (the \emph{compensator}) with $\Lambda(0)=0$ such that
$M(t)=N(t)-\Lambda(t)$ is a local martingale; when $\Lambda$ is absolutely
continuous, $\Lambda(t)=\int_0^t\lambda$.

\begin{intuitionbox}
$\Lambda(t)$ is the predictable ``running forecast'' of how many events should have
happened by $t$; $\lambda$ is its instantaneous rate. Subtract the forecast from
reality and what's left, $M=N-\Lambda$, is pure surprise (a martingale). So
$\lambda$ is exactly the part of the future the past lets you anticipate.
\end{intuitionbox}

\section{The Poisson process and the theorem that runs the book}

\begin{definition}
$N$ is an (inhomogeneous) Poisson process with mean measure $\Lambda$ if for
disjoint sets $A_1,\dots,A_k$ the counts $N(A_i)$ are \emph{independent} with
$N(A_i)\sim\mathrm{Poisson}(\Lambda(A_i))$.
\end{definition}

Two consequences are used everywhere: independent increments, and
$\E[N(t)]=\Lambda(t)$. The converse is the crux.

\begin{theorem}[Watanabe]\label{thm:ch2-watanabe}
A simple point process with a \emph{continuous} compensator $\Lambda$ is a Poisson
process with mean measure $\Lambda$ \emph{if and only if $\Lambda$ is
deterministic}.
\end{theorem}

\begin{whybox}[: this is the whole idea]
Dependence between increments can only come from the intensity \emph{reacting to
history}. If the intensity --- hence the compensator --- is fixed in advance, there
is no channel for the past to influence the future, so the process must be Poisson.
This is precisely why, in Part II, we replace the random Hawkes intensity
$\lambda$ by its deterministic mean $\xi=\E[\lambda]$: doing so \emph{manufactures}
independent increments and hands us a product likelihood.
\end{whybox}

\section{The likelihood (why it needs timestamps)}

\begin{theorem}[Jacod]\label{thm:ch2-jacod}
A point process with intensity $\lambda$ has, on $[0,T]$, log-likelihood (up to a
constant)
\begin{equation}
\ell(\theta)=\sum_{t_i\le T}\log\lambda(t_i;\theta)-\int_0^T\lambda(s;\theta)\dd s.
\label{eq:ch2-ppll}
\end{equation}
\end{theorem}
The first term rewards having forecast high intensity exactly where events landed;
the second penalises total forecast mass. It manifestly needs the times $t_i$ ---
the obstruction of Chapter~\ref{ch:problem}.

\begin{examplebox}[: homogeneous Poisson]
For constant $\lambda$ on $[0,T]$ with $n$ events, \eqref{eq:ch2-ppll} is
$n\log\lambda-\lambda T$; setting the derivative to zero gives the obvious MLE
$\hat\lambda=n/T$. We will reuse this ``rate $=$ count $/$ exposure'' logic when we
estimate the latent immigrant rate from interval-censored data (Ch.~\ref{ch:fitting}).
\end{examplebox}

\begin{recap}
The compensator is the predictable forecast of accumulated events; Watanabe says a
deterministic compensator $\Leftrightarrow$ Poisson $\Rightarrow$ independent
increments; and the point-process likelihood needs timestamps. Part II turns the
first two facts into a method.
\end{recap}

\exercise{Show that for a homogeneous Poisson process the inter-event times are
i.i.d.\ Exponential. (We use the converse --- the time-rescaling theorem --- as a
goodness-of-fit test in Ch.~\ref{ch:identif}.)}
\exercise{Why does Watanabe require the compensator to be \emph{deterministic}, not
merely continuous? Give a process with a continuous but random compensator that is
not Poisson.}
